\begin{lemma}
For every $g,k,f$ as in \noderef{BlockSigma}, the block map $\sigma_f$ is
an increasing injection of $\mathbb N$ into itself.
\end{lemma}
\begin{proof}
Fix $g,k,f$ with $k<g(k)$, and write
\[
G(n)=\noderef{OrbitPoint}(g,k,n)=g^{[n]}(k).
\]
By \noderef{IncSeq}, both $g$ and $f$ are order embeddings of $\mathbb N$
into itself.  In particular, they are strictly increasing.  We first record
the iterate facts needed below.  For every $a$, the map $g^{[a]}$ is strictly
increasing: this is clear for $a=0$, and if it holds for $a$, then
$x<y$ gives
\[
g^{[a+1]}(x)=g\bigl(g^{[a]}(x)\bigr)
 <g\bigl(g^{[a]}(y)\bigr)=g^{[a+1]}(y),
\]
where the middle inequality uses the strict increase of $g$.  Also, induction
on $a$, using the successor rule for iterates, gives
\[
g^{[a]}(g(x))=g^{[a+1]}(x),
\qquad
g^{[a]}\bigl(g^{[b]}(x)\bigr)=g^{[a+b]}(x).
\tag{1}
\]

Applying the first iterate fact to $k<g(k)$ and then using (1) shows
$G(n)<G(n+1)$ for every $n$.  Thus $n<m$ implies $G(n)<G(m)$.  Similarly,
strict increase of $f$ gives
\[
f(n)+1\leq f(n+1).
\tag{2}
\]
An induction using (2) also gives $n\leq f(n)$ for every $n$: the base case
is $0\leq f(0)$, and the successor step follows from
$n+1\leq f(n)+1\leq f(n+1)$.

For each $l$, \noderef{OrbitBlockCoverage} gives either $l<G(0)$ or a
unique $n$ such that
\[
G(n)\leq l<G(n+1).
\]
Here the displayed inequalities are the block condition, with $G$ written
using \noderef{OrbitPoint}.  Define $\sigma(l)$ by the corresponding two
cases:
\[
\sigma(l)=
\begin{cases}
l,&l<G(0),\\
g^{[f(n)-n]}(l),&G(n)\leq l<G(n+1)\text{ for the unique }n.
\end{cases}
\tag{3}
\]
The coverage and uniqueness just proved make (3) a total, well-defined map.
Moreover, (3) is exactly the definition of \noderef{BlockSigma}.  Indeed, in
the second case the chosen $n$ is the least index with $l<G(n+1)$: if $m<n$,
then $m+1\leq n$, so monotonicity of $G$ gives
$G(m+1)\leq G(n)\leq l$.  Thus no earlier index satisfies the defining
inequality.  The coverage lemma ensures that the fallback case in
\noderef{BlockSigma} never occurs.  Consequently
\[
\sigma(l)=\mathsf{BlockSigma}(g,k,f,l)
\qquad\text{for every }l.
\tag{4}
\]

It remains to prove that $\sigma$ is strictly increasing.  On the initial
interval it is the identity, and on each block it is a fixed iterate of $g$;
the iterate fact above proves strict increase within every piece.  For the
first boundary, $G(0)<G(1)$ makes $G(0)$ the first point of the zeroth block,
so (3) and (1) give
\[
\sigma(G(0))=g^{[f(0)]}(G(0))=g^{[f(0)]}(k)=G(f(0)).
\]
Since $0\leq f(0)$ and $G$ is increasing,
\[
l<G(0)\leq G(f(0))=\sigma(G(0))
\tag{5}
\]
for every $l<G(0)$.

Now suppose $G(n)\leq l<G(n+1)$ and set $a=f(n)-n$.  The inequality
$n\leq f(n)$ makes this a natural number.  Strict increase of the iterate,
(1), (2), and the definition of $G$ give
\[
\begin{aligned}
\sigma(l)
 &=g^{[a]}(l)
  <g^{[a]}(G(n+1))\\
 &=g^{[f(n)-n]}\bigl(g^{[n+1]}(k)\bigr)
  =g^{[f(n)+1]}(k)
  =G(f(n)+1)\\
 &\leq G(f(n+1)).
\end{aligned}
\tag{6}
\]
The endpoint $G(n+1)$ belongs to the next block, so (3), (1), and the
inequality $n+1\leq f(n+1)$ give
\[
\begin{aligned}
\sigma(G(n+1))
 &=g^{[f(n+1)-(n+1)]}\bigl(g^{[n+1]}(k)\bigr)\\
 &=g^{[f(n+1)]}(k)=G(f(n+1)).
\end{aligned}
\tag{7}
\]
Combining (6) and (7) proves that every value in one block is strictly less
than the value at the first point of the next block.  Together with (5), the
within-piece inequalities therefore show that $x<y$ implies
$\sigma(x)<\sigma(y)$, including when $x$ and $y$ lie in nonadjacent pieces.
A strictly increasing map on the linearly ordered naturals is an order
embedding, so $\sigma\in\mathsf{IncSeq}$.  Taking this $\sigma$ in (4)
proves the stated existence.
\end{proof}
